% ============================================================
% Close-up: PubMedResearcher — Workflow Interface View
% Fragment — % ============================================================
% Close-up: PubMedResearcher — Workflow Interface View
% Fragment — % ============================================================
% Close-up: PubMedResearcher — Workflow Interface View
% Fragment — % ============================================================
% Close-up: PubMedResearcher — Workflow Interface View
% Fragment — \input{figs/closeup_pubmed}
% ============================================================
\begin{figure}[htbp]
\centering
\resizebox{\textwidth}{!}{%
\begin{tikzpicture}[node distance=0.35cm, font=\sffamily,
    stepbar/.style={rectangle, rounded corners=6pt, thick,
        minimum width=19cm, align=center, font=\sffamily,
        drop shadow={opacity=0.10}, inner sep=6pt, text width=18cm},
    rwpanel/.style={rectangle, rounded corners=5pt, thick,
        minimum width=19cm, align=left, inner sep=6pt,
        font=\sffamily\small, text width=17.5cm},
    infopanel/.style={rectangle, rounded corners=5pt, draw=gray!50, thick,
        minimum width=19cm, align=left, fill=white,
        inner sep=6pt, font=\sffamily\small, text width=17.5cm},
    toolpanel/.style={rectangle, rounded corners=5pt, draw=gray!60, thick,
        minimum width=19cm, align=left, fill=gray!4,
        inner sep=6pt, font=\sffamily\small, text width=17.5cm},
    flowline/.style={->, >=Stealth, very thick, draw=gray!40},
    consumerlabel/.style={font=\sffamily\small\itshape, text=gray!60!black,
        align=center, text width=19cm}
]
    % === STEP BAR ===
    \node[stepbar, draw=violet!50!black, fill=agentPurple] (step)
        {\iconSearch{violet}\;\;
         \textbf{\large PubMedResearcher --- Literature Validation}\\[3pt]
         {\small Planner: \textit{``Literature on TOP targets AND
          pathways --- validates \& classifies novelty.''}}\\[2pt]
         {\scriptsize\sffamily\color{gray!60!black}
          Shown as a third step here; the planner may invoke it
          earlier (e.g.\ for a literature-first query) or re-invoke
          it after replanning.}};

    % === READS ===
    \node[rwpanel, draw=green!40!black, fill=agentGreen,
          below=0.4cm of step] (reads)
        {\textbf{$\blacktriangledown$ Reads from AgentState}\\[5pt]
         $\bullet$\; \texttt{top\_genes} {\scriptsize (list[str])} ---
          ranked DEGs \emph{from OmicMiningAgent};
          prompt says ``search for ALL of them''\\[2pt]
         $\bullet$\; \texttt{query} {\scriptsize (str)} ---
          original research question\\[2pt]
         $\bullet$\; \texttt{previous\_results} {\scriptsize (str)} ---
          accumulated output from Steps 1--2 (omics + KG context)};

    \draw[flowline] (step.south) -- (reads.north);

    % === PROMPT HIGHLIGHTS ===
    \node[infopanel, below=0.4cm of reads] (prompt)
        {\textbf{System Prompt Highlights}
         {\scriptsize (from \texttt{PubMedResearcher.system\_message})}\\[5pt]
         \textbf{1.}\; ``Check \texttt{context["top\_genes"]} for genes
          from previous analysis --- search for ALL of them''\\[2pt]
         \textbf{2.}\; ``Group related genes into 2--4 focused queries
          for efficiency'' {\scriptsize (e.g.\
          \texttt{"TREM2 TYROBP microglia Alzheimer"})}\\[2pt]
         \textbf{3.}\; ``\texttt{pubmed\_full\_search\_tool} is your
          PRIMARY tool --- it downloads and reads full papers''\\[2pt]
         \textbf{4.}\; ``Every claim MUST include a citation:\;
          (Author et al., Year, DOI/PMID)''\\[2pt]
         \textbf{5.}\; Output per gene:\; \textbf{Function},
          \textbf{Disease Role}, \textbf{Key Findings},
          \textbf{Therapeutic Potential}\\[2pt]
         \textbf{6.}\; ``Synthesise findings ACROSS papers ---
          identify convergent and contradictory evidence''};

    \draw[flowline] (reads.south) -- (prompt.north);

    % === TOOLS (two) ===
    \node[toolpanel, below=0.4cm of prompt] (tool)
        {\textbf{Tool 1 (primary):}\;
         \texttt{pubmed\_full\_search\_tool(query, num\_papers=50)}\\[3pt]
         {\scriptsize Downloads full-text PDFs/XMLs from PubMed Central,
          Elsevier, Wiley, arXiv, bioRxiv, medRxiv via paperscraper.\;
          Max 100 papers per call.\;
          Papers cached globally for speed.\;
          Returns metadata + abstract + extracted full text.}\\[6pt]
         \textbf{Tool 2 (fallback):}\;
         \texttt{pubmed\_lite\_tool(query)}\\[3pt]
         {\scriptsize \texttt{[DEPRECATED]} --- lightweight NCBI
          E-utilities search.\; Abstracts and metadata only, no PDF
          download.\; $\leq$20 papers.\; ``Use only when speed is
          critical.''}\\[5pt]
         {\scriptsize \textit{See companion illustration for source
          landscape (PubMed, PMC, Elsevier, Wiley, arXiv, bioRxiv,
          medRxiv).}}};

    \draw[flowline] (prompt.south) -- (tool.north);

    % === WRITES ===
    \node[rwpanel, draw=blue!30!black, fill=stateBlue,
          below=0.4cm of tool] (writes)
        {\textbf{$\blacktriangle$ Writes to AgentState}
         {\scriptsize (via regex extraction from LLM output)}\\[5pt]
         $\bullet$\; \texttt{paper\_dois[]} {\scriptsize (list[str])} ---
          extracted DOI strings for traceability\\[2pt]
         $\bullet$\; \texttt{paper\_pmids[]} {\scriptsize (list[int])} ---
          extracted PubMed ID numbers\\[2pt]
         $\bullet$\; \texttt{agent\_outputs[]} {\scriptsize (str)} ---
          per-gene literature reviews structured as:\\[2pt]
         \hphantom{$\bullet$\;}
          {\scriptsize --- \textbf{Function}: biological role with
           citations}\\[1pt]
         \hphantom{$\bullet$\;}
          {\scriptsize --- \textbf{Disease Role}: mechanism linking
           gene to condition}\\[1pt]
         \hphantom{$\bullet$\;}
          {\scriptsize --- \textbf{Key Findings}: significant results
           across papers}\\[1pt]
         \hphantom{$\bullet$\;}
          {\scriptsize --- \textbf{Therapeutic Potential}: drug targets
           or interventions}\\[1pt]
         \hphantom{$\bullet$\;}
          {\scriptsize --- \textbf{References}: all papers cited with
           DOI/PMID}};

    \draw[flowline] (tool.south) -- (writes.north);

    % === DOWNSTREAM CONSUMERS ===
    \node[consumerlabel, below=0.3cm of writes] (consumers)
        {$\downarrow$\; Consumed by:\;
         \textbf{ScientistsAgent}
          (\texttt{paper\_dois}, literature context)\; $\cdot$\;
         \textbf{ReportingNode}
          (citations, literature-validated targets section)};

\end{tikzpicture}%
}% end resizebox
\caption{\textbf{PubMedResearcher} --- workflow interface view.
Although shown here in the third position, the planner may invoke
this agent at any point---even first for purely literature-driven
queries---and may re-invoke it after replanning if earlier agents
fail or surface unexpected targets.  It provides the evidence layer.
Its system prompt instructs it to read \emph{all} genes from
\texttt{top\_genes} and group them into 2--4 focused queries for
efficient retrieval.  The primary tool downloads up to 100 full-text
papers per call with global caching; a deprecated lite fallback is
available for speed-critical lookups (see companion illustration for
source landscape).  Every claim is required to carry a citation.  The
structured per-gene output and extracted DOIs/PMIDs it writes
to \texttt{AgentState} provide the literature grounding that
ScientistsAgent uses for hypothesis novelty classification and
that ReportingNode uses for the citation apparatus.}
\label{fig:closeup-pubmed}
\end{figure}

% ============================================================
\begin{figure}[htbp]
\centering
\resizebox{\textwidth}{!}{%
\begin{tikzpicture}[node distance=0.35cm, font=\sffamily,
    stepbar/.style={rectangle, rounded corners=6pt, thick,
        minimum width=19cm, align=center, font=\sffamily,
        drop shadow={opacity=0.10}, inner sep=6pt, text width=18cm},
    rwpanel/.style={rectangle, rounded corners=5pt, thick,
        minimum width=19cm, align=left, inner sep=6pt,
        font=\sffamily\small, text width=17.5cm},
    infopanel/.style={rectangle, rounded corners=5pt, draw=gray!50, thick,
        minimum width=19cm, align=left, fill=white,
        inner sep=6pt, font=\sffamily\small, text width=17.5cm},
    toolpanel/.style={rectangle, rounded corners=5pt, draw=gray!60, thick,
        minimum width=19cm, align=left, fill=gray!4,
        inner sep=6pt, font=\sffamily\small, text width=17.5cm},
    flowline/.style={->, >=Stealth, very thick, draw=gray!40},
    consumerlabel/.style={font=\sffamily\small\itshape, text=gray!60!black,
        align=center, text width=19cm}
]
    % === STEP BAR ===
    \node[stepbar, draw=violet!50!black, fill=agentPurple] (step)
        {\iconSearch{violet}\;\;
         \textbf{\large PubMedResearcher --- Literature Validation}\\[3pt]
         {\small Planner: \textit{``Literature on TOP targets AND
          pathways --- validates \& classifies novelty.''}}\\[2pt]
         {\scriptsize\sffamily\color{gray!60!black}
          Shown as a third step here; the planner may invoke it
          earlier (e.g.\ for a literature-first query) or re-invoke
          it after replanning.}};

    % === READS ===
    \node[rwpanel, draw=green!40!black, fill=agentGreen,
          below=0.4cm of step] (reads)
        {\textbf{$\blacktriangledown$ Reads from AgentState}\\[5pt]
         $\bullet$\; \texttt{top\_genes} {\scriptsize (list[str])} ---
          ranked DEGs \emph{from OmicMiningAgent};
          prompt says ``search for ALL of them''\\[2pt]
         $\bullet$\; \texttt{query} {\scriptsize (str)} ---
          original research question\\[2pt]
         $\bullet$\; \texttt{previous\_results} {\scriptsize (str)} ---
          accumulated output from Steps 1--2 (omics + KG context)};

    \draw[flowline] (step.south) -- (reads.north);

    % === PROMPT HIGHLIGHTS ===
    \node[infopanel, below=0.4cm of reads] (prompt)
        {\textbf{System Prompt Highlights}
         {\scriptsize (from \texttt{PubMedResearcher.system\_message})}\\[5pt]
         \textbf{1.}\; ``Check \texttt{context["top\_genes"]} for genes
          from previous analysis --- search for ALL of them''\\[2pt]
         \textbf{2.}\; ``Group related genes into 2--4 focused queries
          for efficiency'' {\scriptsize (e.g.\
          \texttt{"TREM2 TYROBP microglia Alzheimer"})}\\[2pt]
         \textbf{3.}\; ``\texttt{pubmed\_full\_search\_tool} is your
          PRIMARY tool --- it downloads and reads full papers''\\[2pt]
         \textbf{4.}\; ``Every claim MUST include a citation:\;
          (Author et al., Year, DOI/PMID)''\\[2pt]
         \textbf{5.}\; Output per gene:\; \textbf{Function},
          \textbf{Disease Role}, \textbf{Key Findings},
          \textbf{Therapeutic Potential}\\[2pt]
         \textbf{6.}\; ``Synthesise findings ACROSS papers ---
          identify convergent and contradictory evidence''};

    \draw[flowline] (reads.south) -- (prompt.north);

    % === TOOLS (two) ===
    \node[toolpanel, below=0.4cm of prompt] (tool)
        {\textbf{Tool 1 (primary):}\;
         \texttt{pubmed\_full\_search\_tool(query, num\_papers=50)}\\[3pt]
         {\scriptsize Downloads full-text PDFs/XMLs from PubMed Central,
          Elsevier, Wiley, arXiv, bioRxiv, medRxiv via paperscraper.\;
          Max 100 papers per call.\;
          Papers cached globally for speed.\;
          Returns metadata + abstract + extracted full text.}\\[6pt]
         \textbf{Tool 2 (fallback):}\;
         \texttt{pubmed\_lite\_tool(query)}\\[3pt]
         {\scriptsize \texttt{[DEPRECATED]} --- lightweight NCBI
          E-utilities search.\; Abstracts and metadata only, no PDF
          download.\; $\leq$20 papers.\; ``Use only when speed is
          critical.''}\\[5pt]
         {\scriptsize \textit{See companion illustration for source
          landscape (PubMed, PMC, Elsevier, Wiley, arXiv, bioRxiv,
          medRxiv).}}};

    \draw[flowline] (prompt.south) -- (tool.north);

    % === WRITES ===
    \node[rwpanel, draw=blue!30!black, fill=stateBlue,
          below=0.4cm of tool] (writes)
        {\textbf{$\blacktriangle$ Writes to AgentState}
         {\scriptsize (via regex extraction from LLM output)}\\[5pt]
         $\bullet$\; \texttt{paper\_dois[]} {\scriptsize (list[str])} ---
          extracted DOI strings for traceability\\[2pt]
         $\bullet$\; \texttt{paper\_pmids[]} {\scriptsize (list[int])} ---
          extracted PubMed ID numbers\\[2pt]
         $\bullet$\; \texttt{agent\_outputs[]} {\scriptsize (str)} ---
          per-gene literature reviews structured as:\\[2pt]
         \hphantom{$\bullet$\;}
          {\scriptsize --- \textbf{Function}: biological role with
           citations}\\[1pt]
         \hphantom{$\bullet$\;}
          {\scriptsize --- \textbf{Disease Role}: mechanism linking
           gene to condition}\\[1pt]
         \hphantom{$\bullet$\;}
          {\scriptsize --- \textbf{Key Findings}: significant results
           across papers}\\[1pt]
         \hphantom{$\bullet$\;}
          {\scriptsize --- \textbf{Therapeutic Potential}: drug targets
           or interventions}\\[1pt]
         \hphantom{$\bullet$\;}
          {\scriptsize --- \textbf{References}: all papers cited with
           DOI/PMID}};

    \draw[flowline] (tool.south) -- (writes.north);

    % === DOWNSTREAM CONSUMERS ===
    \node[consumerlabel, below=0.3cm of writes] (consumers)
        {$\downarrow$\; Consumed by:\;
         \textbf{ScientistsAgent}
          (\texttt{paper\_dois}, literature context)\; $\cdot$\;
         \textbf{ReportingNode}
          (citations, literature-validated targets section)};

\end{tikzpicture}%
}% end resizebox
\caption{\textbf{PubMedResearcher} --- workflow interface view.
Although shown here in the third position, the planner may invoke
this agent at any point---even first for purely literature-driven
queries---and may re-invoke it after replanning if earlier agents
fail or surface unexpected targets.  It provides the evidence layer.
Its system prompt instructs it to read \emph{all} genes from
\texttt{top\_genes} and group them into 2--4 focused queries for
efficient retrieval.  The primary tool downloads up to 100 full-text
papers per call with global caching; a deprecated lite fallback is
available for speed-critical lookups (see companion illustration for
source landscape).  Every claim is required to carry a citation.  The
structured per-gene output and extracted DOIs/PMIDs it writes
to \texttt{AgentState} provide the literature grounding that
ScientistsAgent uses for hypothesis novelty classification and
that ReportingNode uses for the citation apparatus.}
\label{fig:closeup-pubmed}
\end{figure}

% ============================================================
\begin{figure}[htbp]
\centering
\resizebox{\textwidth}{!}{%
\begin{tikzpicture}[node distance=0.35cm, font=\sffamily,
    stepbar/.style={rectangle, rounded corners=6pt, thick,
        minimum width=19cm, align=center, font=\sffamily,
        drop shadow={opacity=0.10}, inner sep=6pt, text width=18cm},
    rwpanel/.style={rectangle, rounded corners=5pt, thick,
        minimum width=19cm, align=left, inner sep=6pt,
        font=\sffamily\small, text width=17.5cm},
    infopanel/.style={rectangle, rounded corners=5pt, draw=gray!50, thick,
        minimum width=19cm, align=left, fill=white,
        inner sep=6pt, font=\sffamily\small, text width=17.5cm},
    toolpanel/.style={rectangle, rounded corners=5pt, draw=gray!60, thick,
        minimum width=19cm, align=left, fill=gray!4,
        inner sep=6pt, font=\sffamily\small, text width=17.5cm},
    flowline/.style={->, >=Stealth, very thick, draw=gray!40},
    consumerlabel/.style={font=\sffamily\small\itshape, text=gray!60!black,
        align=center, text width=19cm}
]
    % === STEP BAR ===
    \node[stepbar, draw=violet!50!black, fill=agentPurple] (step)
        {\iconSearch{violet}\;\;
         \textbf{\large PubMedResearcher --- Literature Validation}\\[3pt]
         {\small Planner: \textit{``Literature on TOP targets AND
          pathways --- validates \& classifies novelty.''}}\\[2pt]
         {\scriptsize\sffamily\color{gray!60!black}
          Shown as a third step here; the planner may invoke it
          earlier (e.g.\ for a literature-first query) or re-invoke
          it after replanning.}};

    % === READS ===
    \node[rwpanel, draw=green!40!black, fill=agentGreen,
          below=0.4cm of step] (reads)
        {\textbf{$\blacktriangledown$ Reads from AgentState}\\[5pt]
         $\bullet$\; \texttt{top\_genes} {\scriptsize (list[str])} ---
          ranked DEGs \emph{from OmicMiningAgent};
          prompt says ``search for ALL of them''\\[2pt]
         $\bullet$\; \texttt{query} {\scriptsize (str)} ---
          original research question\\[2pt]
         $\bullet$\; \texttt{previous\_results} {\scriptsize (str)} ---
          accumulated output from Steps 1--2 (omics + KG context)};

    \draw[flowline] (step.south) -- (reads.north);

    % === PROMPT HIGHLIGHTS ===
    \node[infopanel, below=0.4cm of reads] (prompt)
        {\textbf{System Prompt Highlights}
         {\scriptsize (from \texttt{PubMedResearcher.system\_message})}\\[5pt]
         \textbf{1.}\; ``Check \texttt{context["top\_genes"]} for genes
          from previous analysis --- search for ALL of them''\\[2pt]
         \textbf{2.}\; ``Group related genes into 2--4 focused queries
          for efficiency'' {\scriptsize (e.g.\
          \texttt{"TREM2 TYROBP microglia Alzheimer"})}\\[2pt]
         \textbf{3.}\; ``\texttt{pubmed\_full\_search\_tool} is your
          PRIMARY tool --- it downloads and reads full papers''\\[2pt]
         \textbf{4.}\; ``Every claim MUST include a citation:\;
          (Author et al., Year, DOI/PMID)''\\[2pt]
         \textbf{5.}\; Output per gene:\; \textbf{Function},
          \textbf{Disease Role}, \textbf{Key Findings},
          \textbf{Therapeutic Potential}\\[2pt]
         \textbf{6.}\; ``Synthesise findings ACROSS papers ---
          identify convergent and contradictory evidence''};

    \draw[flowline] (reads.south) -- (prompt.north);

    % === TOOLS (two) ===
    \node[toolpanel, below=0.4cm of prompt] (tool)
        {\textbf{Tool 1 (primary):}\;
         \texttt{pubmed\_full\_search\_tool(query, num\_papers=50)}\\[3pt]
         {\scriptsize Downloads full-text PDFs/XMLs from PubMed Central,
          Elsevier, Wiley, arXiv, bioRxiv, medRxiv via paperscraper.\;
          Max 100 papers per call.\;
          Papers cached globally for speed.\;
          Returns metadata + abstract + extracted full text.}\\[6pt]
         \textbf{Tool 2 (fallback):}\;
         \texttt{pubmed\_lite\_tool(query)}\\[3pt]
         {\scriptsize \texttt{[DEPRECATED]} --- lightweight NCBI
          E-utilities search.\; Abstracts and metadata only, no PDF
          download.\; $\leq$20 papers.\; ``Use only when speed is
          critical.''}\\[5pt]
         {\scriptsize \textit{See companion illustration for source
          landscape (PubMed, PMC, Elsevier, Wiley, arXiv, bioRxiv,
          medRxiv).}}};

    \draw[flowline] (prompt.south) -- (tool.north);

    % === WRITES ===
    \node[rwpanel, draw=blue!30!black, fill=stateBlue,
          below=0.4cm of tool] (writes)
        {\textbf{$\blacktriangle$ Writes to AgentState}
         {\scriptsize (via regex extraction from LLM output)}\\[5pt]
         $\bullet$\; \texttt{paper\_dois[]} {\scriptsize (list[str])} ---
          extracted DOI strings for traceability\\[2pt]
         $\bullet$\; \texttt{paper\_pmids[]} {\scriptsize (list[int])} ---
          extracted PubMed ID numbers\\[2pt]
         $\bullet$\; \texttt{agent\_outputs[]} {\scriptsize (str)} ---
          per-gene literature reviews structured as:\\[2pt]
         \hphantom{$\bullet$\;}
          {\scriptsize --- \textbf{Function}: biological role with
           citations}\\[1pt]
         \hphantom{$\bullet$\;}
          {\scriptsize --- \textbf{Disease Role}: mechanism linking
           gene to condition}\\[1pt]
         \hphantom{$\bullet$\;}
          {\scriptsize --- \textbf{Key Findings}: significant results
           across papers}\\[1pt]
         \hphantom{$\bullet$\;}
          {\scriptsize --- \textbf{Therapeutic Potential}: drug targets
           or interventions}\\[1pt]
         \hphantom{$\bullet$\;}
          {\scriptsize --- \textbf{References}: all papers cited with
           DOI/PMID}};

    \draw[flowline] (tool.south) -- (writes.north);

    % === DOWNSTREAM CONSUMERS ===
    \node[consumerlabel, below=0.3cm of writes] (consumers)
        {$\downarrow$\; Consumed by:\;
         \textbf{ScientistsAgent}
          (\texttt{paper\_dois}, literature context)\; $\cdot$\;
         \textbf{ReportingNode}
          (citations, literature-validated targets section)};

\end{tikzpicture}%
}% end resizebox
\caption{\textbf{PubMedResearcher} --- workflow interface view.
Although shown here in the third position, the planner may invoke
this agent at any point---even first for purely literature-driven
queries---and may re-invoke it after replanning if earlier agents
fail or surface unexpected targets.  It provides the evidence layer.
Its system prompt instructs it to read \emph{all} genes from
\texttt{top\_genes} and group them into 2--4 focused queries for
efficient retrieval.  The primary tool downloads up to 100 full-text
papers per call with global caching; a deprecated lite fallback is
available for speed-critical lookups (see companion illustration for
source landscape).  Every claim is required to carry a citation.  The
structured per-gene output and extracted DOIs/PMIDs it writes
to \texttt{AgentState} provide the literature grounding that
ScientistsAgent uses for hypothesis novelty classification and
that ReportingNode uses for the citation apparatus.}
\label{fig:closeup-pubmed}
\end{figure}

% ============================================================
\begin{figure}[htbp]
\centering
\resizebox{\textwidth}{!}{%
\begin{tikzpicture}[node distance=0.35cm, font=\sffamily,
    stepbar/.style={rectangle, rounded corners=6pt, thick,
        minimum width=19cm, align=center, font=\sffamily,
        drop shadow={opacity=0.10}, inner sep=6pt, text width=18cm},
    rwpanel/.style={rectangle, rounded corners=5pt, thick,
        minimum width=19cm, align=left, inner sep=6pt,
        font=\sffamily\small, text width=17.5cm},
    infopanel/.style={rectangle, rounded corners=5pt, draw=gray!50, thick,
        minimum width=19cm, align=left, fill=white,
        inner sep=6pt, font=\sffamily\small, text width=17.5cm},
    toolpanel/.style={rectangle, rounded corners=5pt, draw=gray!60, thick,
        minimum width=19cm, align=left, fill=gray!4,
        inner sep=6pt, font=\sffamily\small, text width=17.5cm},
    flowline/.style={->, >=Stealth, very thick, draw=gray!40},
    consumerlabel/.style={font=\sffamily\small\itshape, text=gray!60!black,
        align=center, text width=19cm}
]
    % === STEP BAR ===
    \node[stepbar, draw=violet!50!black, fill=agentPurple] (step)
        {\iconSearch{violet}\;\;
         \textbf{\large PubMedResearcher --- Literature Validation}\\[3pt]
         {\small Planner: \textit{``Literature on TOP targets AND
          pathways --- validates \& classifies novelty.''}}\\[2pt]
         {\scriptsize\sffamily\color{gray!60!black}
          Shown as a third step here; the planner may invoke it
          earlier (e.g.\ for a literature-first query) or re-invoke
          it after replanning.}};

    % === READS ===
    \node[rwpanel, draw=green!40!black, fill=agentGreen,
          below=0.4cm of step] (reads)
        {\textbf{$\blacktriangledown$ Reads from AgentState}\\[5pt]
         $\bullet$\; \texttt{top\_genes} {\scriptsize (list[str])} ---
          ranked DEGs \emph{from OmicMiningAgent};
          prompt says ``search for ALL of them''\\[2pt]
         $\bullet$\; \texttt{query} {\scriptsize (str)} ---
          original research question\\[2pt]
         $\bullet$\; \texttt{previous\_results} {\scriptsize (str)} ---
          accumulated output from Steps 1--2 (omics + KG context)};

    \draw[flowline] (step.south) -- (reads.north);

    % === PROMPT HIGHLIGHTS ===
    \node[infopanel, below=0.4cm of reads] (prompt)
        {\textbf{System Prompt Highlights}
         {\scriptsize (from \texttt{PubMedResearcher.system\_message})}\\[5pt]
         \textbf{1.}\; ``Check \texttt{context["top\_genes"]} for genes
          from previous analysis --- search for ALL of them''\\[2pt]
         \textbf{2.}\; ``Group related genes into 2--4 focused queries
          for efficiency'' {\scriptsize (e.g.\
          \texttt{"TREM2 TYROBP microglia Alzheimer"})}\\[2pt]
         \textbf{3.}\; ``\texttt{pubmed\_full\_search\_tool} is your
          PRIMARY tool --- it downloads and reads full papers''\\[2pt]
         \textbf{4.}\; ``Every claim MUST include a citation:\;
          (Author et al., Year, DOI/PMID)''\\[2pt]
         \textbf{5.}\; Output per gene:\; \textbf{Function},
          \textbf{Disease Role}, \textbf{Key Findings},
          \textbf{Therapeutic Potential}\\[2pt]
         \textbf{6.}\; ``Synthesise findings ACROSS papers ---
          identify convergent and contradictory evidence''};

    \draw[flowline] (reads.south) -- (prompt.north);

    % === TOOLS (two) ===
    \node[toolpanel, below=0.4cm of prompt] (tool)
        {\textbf{Tool 1 (primary):}\;
         \texttt{pubmed\_full\_search\_tool(query, num\_papers=50)}\\[3pt]
         {\scriptsize Downloads full-text PDFs/XMLs from PubMed Central,
          Elsevier, Wiley, arXiv, bioRxiv, medRxiv via paperscraper.\;
          Max 100 papers per call.\;
          Papers cached globally for speed.\;
          Returns metadata + abstract + extracted full text.}\\[6pt]
         \textbf{Tool 2 (fallback):}\;
         \texttt{pubmed\_lite\_tool(query)}\\[3pt]
         {\scriptsize \texttt{[DEPRECATED]} --- lightweight NCBI
          E-utilities search.\; Abstracts and metadata only, no PDF
          download.\; $\leq$20 papers.\; ``Use only when speed is
          critical.''}\\[5pt]
         {\scriptsize \textit{See companion illustration for source
          landscape (PubMed, PMC, Elsevier, Wiley, arXiv, bioRxiv,
          medRxiv).}}};

    \draw[flowline] (prompt.south) -- (tool.north);

    % === WRITES ===
    \node[rwpanel, draw=blue!30!black, fill=stateBlue,
          below=0.4cm of tool] (writes)
        {\textbf{$\blacktriangle$ Writes to AgentState}
         {\scriptsize (via regex extraction from LLM output)}\\[5pt]
         $\bullet$\; \texttt{paper\_dois[]} {\scriptsize (list[str])} ---
          extracted DOI strings for traceability\\[2pt]
         $\bullet$\; \texttt{paper\_pmids[]} {\scriptsize (list[int])} ---
          extracted PubMed ID numbers\\[2pt]
         $\bullet$\; \texttt{agent\_outputs[]} {\scriptsize (str)} ---
          per-gene literature reviews structured as:\\[2pt]
         \hphantom{$\bullet$\;}
          {\scriptsize --- \textbf{Function}: biological role with
           citations}\\[1pt]
         \hphantom{$\bullet$\;}
          {\scriptsize --- \textbf{Disease Role}: mechanism linking
           gene to condition}\\[1pt]
         \hphantom{$\bullet$\;}
          {\scriptsize --- \textbf{Key Findings}: significant results
           across papers}\\[1pt]
         \hphantom{$\bullet$\;}
          {\scriptsize --- \textbf{Therapeutic Potential}: drug targets
           or interventions}\\[1pt]
         \hphantom{$\bullet$\;}
          {\scriptsize --- \textbf{References}: all papers cited with
           DOI/PMID}};

    \draw[flowline] (tool.south) -- (writes.north);

    % === DOWNSTREAM CONSUMERS ===
    \node[consumerlabel, below=0.3cm of writes] (consumers)
        {$\downarrow$\; Consumed by:\;
         \textbf{ScientistsAgent}
          (\texttt{paper\_dois}, literature context)\; $\cdot$\;
         \textbf{ReportingNode}
          (citations, literature-validated targets section)};

\end{tikzpicture}%
}% end resizebox
\caption{\textbf{PubMedResearcher} --- workflow interface view.
Although shown here in the third position, the planner may invoke
this agent at any point---even first for purely literature-driven
queries---and may re-invoke it after replanning if earlier agents
fail or surface unexpected targets.  It provides the evidence layer.
Its system prompt instructs it to read \emph{all} genes from
\texttt{top\_genes} and group them into 2--4 focused queries for
efficient retrieval.  The primary tool downloads up to 100 full-text
papers per call with global caching; a deprecated lite fallback is
available for speed-critical lookups (see companion illustration for
source landscape).  Every claim is required to carry a citation.  The
structured per-gene output and extracted DOIs/PMIDs it writes
to \texttt{AgentState} provide the literature grounding that
ScientistsAgent uses for hypothesis novelty classification and
that ReportingNode uses for the citation apparatus.}
\label{fig:closeup-pubmed}
\end{figure}
